Five questions our replication cannot answer with the artefacts on disk, listed with their basis in what \emph{was} run vs.\ what \emph{was not}, and with a concrete next-step probe each.

\begin{enumerate}

\item \textbf{Does the proven prefactor advantage of Nakaji-FAE over Grinko-IQAE ($4.1\times 10^{3}$ vs $1.15\times 10^{6}$, $\sim\!280\times$) survive a head-to-head empirical race on identical $(a,\epsilon)$ targets, and does the empirical prefactor gap track the ratio of proven bounds or collapse in practice?}
\emph{Basis.} Our replication only sanity-checked FAE's own empirical prefactor against its own upper bound; we did \emph{not} implement Grinko IQAE and race the two on matched amplitudes. This is the single most publishable single-number comparison in the paper and the largest gap in the current replication (Critique item~1).
\emph{Next steps.} Implement Grinko IQAE (arXiv:1912.05559) as a third algorithm in \texttt{code/iqae.py}, reusing the existing \texttt{oracle.py} Grover operator; run the same \texttt{experiment.py} sweep for $a\in\{0.1,0.2,0.3,0.4\}$ at matched $N_{\mathrm{orac}}$ budgets; extract per-$a$ prefactors from the $\log_{10}(N_{\mathrm{orac}}) = 1\cdot\log_{10}(1/\epsilon) + \log_{10} C$ fits (forced Heisenberg slope) and report the empirical FAE-to-IQAE prefactor ratio alongside the proven-bound ratio; if the empirical ratio is much smaller than 280, quantify the effective advantage.

\item \textbf{Does FAE's two-stage design retain its edge over both IQAE and MLAE under realistic depolarising and readout noise, or does the Chernoff-derived shot count ($N_{\mathrm{shot}}^{1\mathrm{st}} = 1944\ln(2/\delta_{c})$) get eaten by noise-induced bias in the $\cos(2(2^{k}+1)\theta)$ estimates?}
\emph{Basis.} Everything in this replication is noiseless statevector (Critique item~3). FAE's marketing pitch is NISQ-friendliness, so its behavior under realistic noise \emph{is} the point --- and we did not test it. A biased cosine estimate would bleed into $\arccos$ in stage~1 and into the trig-addition in stage~2; the failure mode is not obvious.
\emph{Next steps.} Wrap \texttt{oracle.py}'s \texttt{exact\_prob\_good\_after\_Qm} with a depolarising channel $p_{\mathrm{depol}} \in \{10^{-4}, 10^{-3}, 10^{-2}\}$ per $Q^{m}$ application (composition of $\sim 4m$ gate errors) and a readout confusion matrix $(p_{01},p_{10}) \in \{0, 0.01, 0.03\}$; rerun FAE, IQAE, and MLAE at $a=0.3$ across the noise grid; report which algorithm's slope stays closest to~1 and at what noise level FAE loses its rigorous-guarantee regime (where the Chernoff bound becomes vacuous).

\item \textbf{Can a noise-adaptive query schedule (increase $N_{\mathrm{shot}}$ per round when the observed cosine variance exceeds the Chernoff-implied variance) recover FAE's Heisenberg scaling under noise more efficiently than fixed conservative shot counts?}
\emph{Basis.} FAE's $N_{\mathrm{shot}}^{1\mathrm{st}} = 1944\ln(2/\delta_{c})$ is set to satisfy the worst-case Chernoff bound (Critique item~2, $\sim\!10\,300$ shots at $\delta_{c}=0.01$). Under noise this is either under-shot (biased) or over-shot (wasted queries); an adaptive schedule trades rigor for query economy.
\emph{Next steps.} Implement \texttt{fae\_adaptive.py}: after each \texttt{COS(m, N\_shot\_min=100)} call, compute the empirical variance of $c_{m}$ over a small pilot batch, compare to the Chernoff-implied $1/(2 N_{\mathrm{shot,min}})$, and if the pilot variance is $>\!2\times$ expected, double $N_{\mathrm{shot}}$ until the ratio is $\leq\!2$; measure $N_{\mathrm{orac}}(\epsilon)$ under the same noise grid as Q2 and compare to fixed-shot FAE; report the query-economy ratio at $\epsilon = 10^{-3}$.

\item \textbf{Does replacing FAE's iteration schedule $m_{k} = 2^{k}$ with an ML-selected schedule (a small policy network trained to pick the next $m_{k}$ given $[\theta_{\min}, \theta_{\max}]$ and shot budget) reduce the per-$a$ prefactor $C$ without breaking the Heisenberg $1/\epsilon$ scaling?}
\emph{Basis.} The $2^{k}$ doubling schedule is a natural, algorithm-agnostic default (also used by IQAE and MLAE) but nothing forces it to be optimal for a given $a$; our fits give per-$a$ prefactors $C \approx 88.9$ at $a=0.3$ with room to shave. Candidate probe explicitly listed by Rick.
\emph{Next steps.} Frame schedule selection as a small MDP: state $(\theta_{\min}, \theta_{\max}, k, N_{\mathrm{budget}})$, action $m_{k+1}$ in a bounded set, reward $-\log|\hat a - a|$; train a policy on offline traces from \texttt{sweep\_raw.csv} via behavior cloning + REINFORCE with a KL penalty to the $2^{k}$ baseline; evaluate on held-out $a$ values and report the median $N_{\mathrm{orac}}$ reduction at fixed $\epsilon$.

\item \textbf{Does FAE deliver a measurable end-to-end speedup for a canonical quant-finance application (European call-option pricing via QAE, Woerner--Egger 2019, arXiv:1806.06893) at commercially interesting error targets ($\epsilon \in [10^{-3}, 10^{-4}]$), and is the constant-factor advantage over IQAE preserved once the payoff-loading circuit is included in the query count?}
\emph{Basis.} AE's most-cited application is quant-finance Monte-Carlo replacement; the paper positions FAE as the NISQ-friendly successor. If the payoff-loading oracle $A_{\mathrm{payoff}}$ dominates the query count (common for realistic loss distributions), FAE's constant-factor win on the $Q$-application count may or may not translate to end-to-end wallclock. Candidate probe explicitly listed by Rick.
\emph{Next steps.} Reuse Qiskit-Finance's \texttt{EuropeanCallPricing} (or reimplement its $A_{\mathrm{load}}$ + payoff comparator as a standalone module) as the $A$ oracle in \texttt{oracle.py}; run FAE, IQAE, and canonical QPE-QAE (small-qubit) at $\epsilon \in \{10^{-2}, 10^{-3}\}$ for strike-price grids typical of Woerner--Egger; report total oracle calls (loading + $Q$) per $\epsilon$ and end-to-end call-price RMSE; flag whether FAE's algorithm-level advantage over IQAE survives the loading-cost dilution.

\end{enumerate}
